\documentclass[12pt,aspectratio=169]{beamer}  
%\usetheme{Madrid}
%\usetheme{default}                 
\setbeamertemplate{navigation symbols}{}
\setbeamertemplate{footline}{\color{gray}\hrule \hfill \scriptsize Corso di Economia  e Finanza Computazionale\insertslidenavigationsymbol \insertframenavigationsymbol \scriptsize \insertframenumber/\inserttotalframenumber \hspace*{2mm}}
\setbeamertemplate{headline}{\color{gray}\vskip1mm \hskip2mm \scriptsize UNICH, LM56 CLEF - Economia e Finanza \vskip0.5mm \hrule}
%\setbeamertemplate{headline}{\hskip0.5mm\includegraphics[scale=0.2]{../../../../loghi/logolm56.png}}
%\usepackage[italian]{babel}
%\usepackage[T1]{fontenc}
\usepackage[utf8]{inputenc}
\usepackage{longtable,booktabs}

\begin{document}
\begin{frame}
\unitlength 5mm
\begin{picture}(25,1)
\linethickness{2pt}
\put(2,0.5){\line(1,0){20}}
\end{picture}
\begin{center}
ECONOMIA E FINANZA COMPUTAZIONALE

La struttura a termine dei tassi di interesse
\end{center}

\begin{picture}(25,1)
\linethickness{2pt}
\put(2,0.5){\line(1,0){20}}
\end{picture}
\end{frame}

\input{"slides.tex"}

\end{document}
